\section{Conclusion and future scope}
\label{sec:conclusion}

RoomRecon turns an iPhone capture of a room into a photoreal Gaussian splat that can be walked through with collision, using maintained tools for almost every stage and our own code for the glue and the evaluation. Taking poses and depth from the phone, instead of estimating poses from images as in our first attempt, removed the failure on plain walls, and DN-Splatter's priors made the geometry 3.3 times more accurate than plain splats (1.76 vs 5.85\,cm from the laser scan) for 0.5\,dB less PSNR near the training walk. A public room and our own room were both walkable on a 6\,GB laptop GPU. Our main measurement shows that walk collision built from DN-Splatter's mesh matches a laser scan far better than the default collision built from the splat: the walkable-area IoU rose from 0.58 to 0.81 and missed walls fell from 17.8\% to 1.7\%, at the cost of 0.04\,m$^2$ of phantom obstacles, which failed one of our four pre-agreed criteria. Polycam's mesh of the same capture performed almost as well, both as collision (0.79, 1.8\%) and as geometry (1.86\,cm from the scan), so the main gain is over the default splat collision, not over existing apps. The gap was not felt in a free walk and the study covers one room, so we present it as a measured difference, not a general law. Two other novelty attempts, a capture check and a tracking repair, failed the tests we set in advance; we report them because they show what did not work and why.

Two directions follow from this work:
\begin{enumerate}
\item \textbf{Repeat the collision study on more rooms.} MuSHRoom provides reference meshes for all ten of its rooms, which would show whether the advantage of mesh-based collision holds beyond \texttt{vr\_room} and whether the mesh should become the default collision source.
\item \textbf{Poses from images for phones without LiDAR.} Feed-forward models such as VGGT \cite{wang2025vggt} and MASt3R-SfM \cite{duisterhof2025mast3rsfm} could be tested on the plain-walled bedroom that defeated our v1.
\end{enumerate}

\section*{Acknowledgements}
We thank IIIT Kottayam for access to its GPU server, and Kaggle for free GPU sessions. We used the MuSHRoom dataset \cite{ren2024mushroom} (CC BY 4.0) and, in v1, Apple's ARKitScenes \cite{baruch2021arkitscenes}. Our pipeline is built on open-source work by others: nerfstudio \cite{tancik2023nerfstudio}, DN-Splatter \cite{turkulainen2025dnsplatter, dnsplattercode}, gsplat \cite{ye2025gsplat}, Omnidata \cite{eftekhar2021omnidata}, Open3D \cite{zhou2018open3d}, PlayCanvas' splat-transform \cite{splattransform} and SuperSplat Viewer \cite{supersplatviewer}, and pixi \cite{pixi}; captures were made with Polycam \cite{polycam}. In v1 we also used COLMAP \cite{schoenberger2016colmap}, hloc \cite{sarlin2019hloc} with SuperPoint \cite{detone2018superpoint} and LightGlue \cite{lindenberger2023lightglue}, ORB-SLAM3 \cite{campos2021orbslam3}, GaussianSplats3D \cite{gaussiansplats3d} and NeRFCapture \cite{nerfcapture}.
